\section{An exact comparison of the two regulators}
\label{sec:rg}

The spectral calculation concerns the Brownian regulator, whereas the
mass in the theorem belongs to the Wilson measure.  We connect them by
comparing their partition functions, including a central twist, on
rectangles whose lengths grow with the coupling.  The required error is
exponentially small in $\beta$ relative to the untwisted partition
function.  This is stronger than equality of limiting pressures: it
will resolve an odd-sector trace which is itself exponentially small in
the temporal length.

The construction builds on the exact blocking transformation of
\cite[Sections ``The linear block transformation'', ``An exact
renormalization map'', and ``Admission and terminal-coupling
matching'']{SharpO4}.  We state the parts of that construction used below
and prove the extensions needed here.  There are four changes: independent
spatial and temporal stiffnesses, a real lift of a circle coordinate,
complex local weights, and integer curvature.  In particular, the
positive-density comparison theorem of that reference cannot directly
compare our two regulators.

\subsection{The comparison to be established}

In this and the next two sections, $L$ denotes a fixed large dyadic
blocking factor and $N$ denotes blocking depth.  For a prescribed large
Wilson coupling $\beta$, the admission and shooting proposition of
\cite[Section ``Admission and terminal-coupling matching'']{SharpO4}
selects a depth $N$ with
\[
                    L^N\asymp e^{\pi\beta}/\sqrt\beta.
\]
The corresponding Wilson trajectory lies inside its admitted bands.
Write a rectangle's fine periods as $L^N m_x,L^N m_t$; they are even
integers.  The subscripts $+$ and $-$ will denote periodic and
minus-center conditions in time, respectively.

We augment the Wilson measure by a Gaussian \emph{flat lift}.  Let $v$
have period $\pi$.  On the universal cover of the periodic rectangle,
choose real values of $v$ with increments $\pi n_i$ on traversing period
$i$, where $(n_x,n_t)\in\mathbb Z^2$.  Integrate one value over a
fundamental interval and the remaining relative values over $\mathbb R$,
sum over the two winding integers, and use the energy
\[
                   \frac12\sum_{i=x,t}\sum_{e_i}c_i(d_i v)^2,
                   \qquad c_i=\beta+O(1).
\]
In that summand give the Wilson field the additional central twist
$(-1)^{n_i}$ in direction $i$.  Denote the resulting partition function
by $Z^{F(0)}$.  The superscript $F(\phi)$ means that the integrand is
also multiplied by $\exp\{i\phi\sum_{e_t}d_t v\}$, with true increments
on the cover.  A change of the one-site normalization of this integral
only changes a volume scalar.

\begin{proposition}[Regulator comparison]\label{prop:rg-comparison}
For all sufficiently large $\beta$ there are
\[
 b=\beta+O(1),\qquad r=1+O(\beta^{-1}),\qquad
 c_x,c_t=\beta+O(1),\qquad \phi\in\mathbb R,
\]
and a bulk volume scalar for each regulator, independent of the periods
and of the time seam, with the following property.  Remove these scalars
from the partition functions.  For dyadic lengths
\[
 m_x\asymp\beta(\log\beta)^2,
 \qquad m_t\asymp\sqrt\beta(\log\beta)^4,
\]
and also with $m_t$ replaced by $2m_t$, the Brownian regulator of heat
time $r/b$ per fine horizontal cell, hence circumference
$W=(r/b)L^Nm_x$, satisfies
\begin{equation}
 Z^B_\pm=Z^{F(\phi)}_\pm+
                     O(e^{-c\beta})Z^{F(0)}_+ .
 \label{eq:R1}
\end{equation}
Here $c>0$ and the implicit constants are independent of $\beta$ and
these periods.  The same scalar for a given regulator is used for both
seams and both temporal lengths.
\end{proposition}

We prove the exact blocking estimates in this section.
Section~\ref{sec:initialization} places both microscopic regulators in
the resulting class.  Section~\ref{sec:comparison} chooses the four
matching parameters and integrates the terminal discrepancy, completing
the proof of Proposition~\ref{prop:rg-comparison}.  The finite Taylor
coefficient maps are retained throughout, since they also determine the
additive coupling renormalization in Section~\ref{sec:renormalization}.

\subsection{Scales, lifts, and the Gaussian operators}

Set
\begin{equation}
 \begin{gathered}
 \gamma=\pi^{-1}\log L,\qquad H_j=H+\gamma j,\qquad
 M_j=\lceil(\log H_j)^2\rceil,\\
 p_j=(\log H_j)^{P_0},\qquad t_j=H_j^{-1/2}p_j,
 \qquad T_j=H_j^{-.49},\qquad h=N-j.
 \end{gathered}
 \label{eq:rg-scales}
\end{equation}
The running layers are $j=N,N-1,\ldots,0$.  The fixed power $P_0$ will
be larger than $20$ and than finitely many powers arising below; $H$ is
chosen last.  A layer has periods $m_xL^j,m_tL^j$, with
$\min(m_x,m_t)\ge C_{\rm adm}M_0$.  All cutoffs use the reference value
$H_j$, not a running kinetic coefficient.  This makes the cutoffs common
to two trajectories being compared.

At layer $j$, put
\[
 R_j=\sqrt{\beta/H_j},\qquad P_j=\pi R_j,\qquad w=R_jv.
\]
At a site integrate $(q,w)\in S^3\times[0,P_j)$ with measure
$d\omega_3(q)\,dw$.  Assign an integer $k_e$ to each oriented positive
lattice edge and set
\[
 u_e=d_ew+P_jk_e,
 \qquad \nabla_eq=q_y-(-1)^{k_e}\varepsilon_eq_x,
 \qquad r_e^2=|\nabla_eq|^2+u_e^2.
\]
Here $\varepsilon_e=-1$ on the specified central seam and is $1$
elsewhere.  Replacing a site representative by $w_x+P_jn_x$,
$q_x\mapsto(-1)^{n_x}q_x$, and changing $k$ to keep $u$ fixed is a change
of coordinates.  All expressions below are invariant under it.  The
flat-lift reference is exactly the restriction $dk=0$ on every
plaquette; the Brownian density permits $dk\ne0$.  The lower bound on
$R_j$ is uniform, although its upper bound need not be.

There are four displayed kinetic coefficients,
$B_x^q,B_t^q,B_x^w,B_t^w$.  We use
\[
 B=(B_x^qB_t^q)^{1/2},\qquad
 \chi=(B_x^wB_t^w)^{1/2}/B,
\]
and write their determinant-one shapes as $(l_q^{-1},l_q)$ and
$(l_w^{-1},l_w)$.  The shapes lie in a fixed small neighborhood of one,
$\chi\in[1/2,2]$, and $B\in[H_j/2,2H_j]$.  Slightly enlarged bands are
allowed during a single step.  Symmetries are left and right quaternion
multiplication, common $w$ translation, spatial reflection, time
reflection accompanied by $w\mapsto-w$, and $w\mapsto-w$ accompanied by
complex conjugation.  They imply that an affine quadratic form has
exactly these four coefficients.  A mixed $q$--$w$ bilinear is excluded
by quaternion conjugation symmetry, and spatial reflection excludes
mixed spatial--temporal derivatives.

Let $Q$ be the block average with weights obtained by sampling and
normalizing the product of an even smooth bump supported in the middle
of an $L\times L$ block.  Thus $Q1=1$ and its first moment is the block
center.  Observation precision is one.  We distinguish the conditional
mean operator $\mathsf P_h$ below from the lift period $P_j$.

For the flat reference the initial scalar precision of shape $l$ is the
nearest-neighbor precision $l^{-1}D_x+lD_t$.  For the Brownian regulator
it is the Schur precision obtained by minimizing, in each horizontal
strip,
\[
 \frac12l^{-1}\int_0^1\sum_t\dot f_t^2\,ds
       +\frac12l\int_0^1\sum_t|f_{t+1}-f_t|^2\,ds
\]
with prescribed endpoint columns.  The following explicit edge lift
will also be used in the initialization.

\begin{lemma}[Free operators for the two initial precisions]
\label{lem:rg-free}
For either initial precision and either admissible shape, the free block
iteration has uniformly elliptic, exponentially localized precisions
$K_h$, conditional means $\mathsf P_h$, conditional covariances $C_h$,
and positive edge stacks.  Their required fixed derivatives, including
shape derivatives, obey bounds on one common complex momentum
neighborhood.  For two initial precisions with the same shape, their
histories differ by $CL^{-h}$ in those bounds.  Constants needed to
choose $L$ are independent of sufficiently large $L$.

For a Brownian strip the edge lift consists of the horizontal term
$l^{-1}$ and a positive matrix on the two endpoint values of each
vertical edge.  If $z$ is the vertical edge Laplacian symbol, this matrix
is
\begin{equation}
 B_l(z)=l\sum_{m\ge1}J_m^*
          \bigl(1+l^2z/(\pi m)^2\bigr)^{-1}J_m,
 \qquad J_s(c,c')=(1-s)c+sc',
 \label{eq:R3}
\end{equation}
where $J_m$ is the $m$th normalized sine coefficient of $J_s$.  It is
analytic near $[0,4]$, with uniform positive lower and upper bounds on
that real interval.  The formula holds also for cut vertical arrays,
using their existing bonds.  The affine precision, after strips are
summed, is $\operatorname{diag}(l^{-1},l)$.
\end{lemma}

\begin{proof}
Subtract $J_s$ from a strip field and expand its zero-endpoint part in
sine modes.  Completing the square at eigenvalue $(\pi m)^2$ gives
\eqref{eq:R3}.  At $z=0$ its quadratic form is
$l\int_0^1|J_s(c,c')|^2ds$, which is positive for a nonzero endpoint
pair.  The resolvent is positive on $[0,4]$; compactness gives uniform
bounds, and the poles stay away from a common neighborhood of that
interval.  Summing strips gives the stated affine form.

We recall the quantitative free input from
\cite[Section ``The linear block transformation'' and Appendix
``Uniform analytic estimates for the linear block map'', Lemmas
``Uniform shells and strips'' and ``Alias syzygy and coarse
division'']{SharpO4}.  If $b_Q(k)$ is the Fourier symbol of $Q$, then
\[
 K_{h+1}(p)^{-1}=1+L^{-2}\sum_{l\bmod L}
      b_Q(k_l)b_Q(k_l)^*K_h(k_l)^{-1},\qquad
 k_l=(p+2\pi l)/L.
\]
With stars continued analytically from real momentum, the explicit
depth formula has principal numerator $m^2K_0(p/m)$, where $m=L^h$.
The product of bump symbols has nonprincipal shell sums bounded by
$CA_0^h$.  The only estimates on the initial numerator required in
that proof are even analyticity on a common strip, bounded derivatives,
real ellipticity, and its specified quadratic jet with a fourth-order
remainder.  Both of our precisions have these properties.  In
particular, on every nonprincipal alias $p_l=p+2\pi l$,
\[
 |m^2K_0(p_l/m)|\ge c(1+|l|)^2,
 \qquad
 |m^2(K_{0,1}-K_{0,2})(p_l/m)|
                         \le C|p_l|^4/m^2.
\]
Real ellipticity proves the first estimate on the real cube; reducing
the imaginary width preserves it.  Even Taylor expansion on a larger
analytic cube proves the second.  The reciprocal difference is therefore
$O(m^{-2})$.  Summation with the product-bump weights costs
$CA_0^h$, and $L>A_0$ gives $CL^{-h}$.  The principal alias has its direct
Taylor estimate.  The principal denominator stays bounded away from
zero on a common smaller strip, so Cauchy estimates give every fixed
derivative, including the shape derivatives.

For clarity, the other free conclusions we use are
\begin{align}
 &C_h=(K_h+Q^*Q)^{-1},\qquad
 K_h\mathsf P_h=Q^*K_{h+1},\qquad I-Q\mathsf P_h=K_{h+1},
 \nonumber\\
 &\mathsf P_h1=1,\qquad
   \mathsf P_h(\text{centered affine field})
                         =\text{its affine extension},\nonumber\\
 &\sup_x\sum_y e^{c|x/L-y|}
   |\nabla_f^s\mathsf P_h(x,y)|\le C_sL^{-s},\qquad s=0,1,2.
 \label{eq:rg-free-prediction}
\end{align}
The covariance has exponentially decaying kernels in block units, with
constants allowed to depend on fixed $L$.

Here is the edge form of the positive completion imported from
\cite[Lemma ``Positive completion of the block transformation'']{SharpO4}.
Choose a common $c_0>0$ and write
$\mathsf A_h=(\sqrt{c_0}I,\ell_h)^\mathsf t$, so that
$K_h=d^*\mathsf A_h^*\mathsf A_hd$.  There are operators
$\mathcal M_h,\mathcal N_h$ with
\begin{align}
 \mathcal M_h^*\mathcal M_h+\mathcal N_h^*\mathcal N_h&=I,\nonumber\\
 \mathcal M_h^*(\mathsf A_hdf,V-Qf)&=\mathsf A_{h+1}d'V,\nonumber\\
 \mathcal M_h\mathsf A_{h+1}d'
       &=(\mathsf A_hd\mathsf P_h,I-Q\mathsf P_h),\qquad
 \mathcal N_h\mathsf A_{h+1}d'=0.
 \label{eq:rg-positive-completion}
\end{align}
The part of $\mathcal M_h$ taking coarse data to fine edge rows has
weighted row bound $C/L$ and column bound $CL$; its first two fine
differences have bounds $C/L^2,C/L^3$.  The coarse parts have bounded
row and column sums.  All these statements retain the history gain.

For our shapes, start the raw edge lift at the shape diagonal rather
than the identity.  Division of its fourth-order-zero remainder and
the curl addition used in the cited proof remain uniformly bounded.
The two principal linear jets in the construction of the edge
intertwiner have the same shape factor; hence both alias divisions
and their gluing apply.  The common real bounds give positivity.
At depth zero retain the specified strip lift without adding the fine
inverse-curl correction; deleting that correction improves the
positive-defect inequality in the cited proof.  Its second-stack rows
are the square root after removing $c_0/2$ on each side of
\eqref{eq:R3}.  Zero rows can be added to give common index sets.
Finally take the direct sum of the $q$ and $w$ constructions, multiplying
all $w$ edge, observation, and propagated output entries by
$\sqrt\chi$.  This proves the assertions for the four kinetic
coefficients.
\end{proof}

\subsection{Densities retained by the transformation}

We now specify what is controlled after a block integration.  In a region
where all increments are small, finitely many polynomial interactions
are displayed and the remaining function has a derivative norm.  A
region containing large increments is represented instead by a bounded
weight with an exact inventory of those increments.  Recording every
argument read by such a weight is essential to the later product
estimates.

The actual energy uses \emph{numerical} second-stack rows: truncate their
stencils at radius $c_LM_j$, with $c_L$ small enough for all compositions
in \cite[Section ``The energy identity and the massive Gaussian'' and
Appendix ``Local sector factors and their complete supports'']{SharpO4}.
For \eqref{eq:R3}, truncate a Chebyshev expansion of the square root.
Analyticity gives error $C_Le^{-c_LM_j}$ in operator norm, also for cut
arrays; bandwidth per polynomial degree then gives weighted row and
column errors.  Correct the same-center coefficient, by the same
exponential amount, so that affine normalization is exact.  Coefficient
calculations always use full kernels.  Numerical rows are needed in the
actual density because a finite-read mask cannot bound an infinite row
acting on unbounded $u$.

A $q$ stencil is transported along specified coordinate paths, averaged
over the reflection-related choices.  Its complete read set includes
those paths.  A second-stack row centered at $e$ is enabled only if
\[
 r_f\le t_j\exp\{\operatorname{dist}(e,f)/R_*\}
\]
for every edge on its numerical stencil and its transport paths, and if
there is no edge with $r_f>t_j$ within fixed distance $D_0$ of its
center.  The fixed $R_*$ is larger than the localization length.  The
initialization will choose $D_0$ before $L$.  Enabled rows contribute
half their squared norm to $\mathcal E$.  Direct rows contribute
\begin{equation}
 S_t(e)=
 \begin{cases}
 \dfrac{c_0}{2}\bigl(|\nabla_eq|^2+\chi u_e^2\bigr),&r_e\le t,\\[2mm]
 d_0t\dfrac{|\nabla_eq|^2+\chi u_e^2}{r_e},&r_e>t,
 \end{cases}
 \qquad 0<d_0<c_0/16,
 \label{eq:R4}
\end{equation}
where $t=t_j$.  The tail is linear in $r_e$ and remains integrable in
the lifted variable.

Here are the locality conventions used for every label.  A label has an
anchor, its complete read support, and a connected record of cubes of
radius $M_j$ and joining paths.  Its integer load $s\ge1$ bounds the
number of sites by $CM_j^2s$ and the record length by $CM_js$.
Every spin argument, path, mask, eligibility test, negative status query,
and covariance window is recorded.  A finite-period label retains its
lifted record before folding.  Intersecting records can be joined at a
fixed multiple of total load.  Projection to the next lattice, including
the fixed enlargements used below, has load
\begin{equation}
               4s/L\le s'\le D_*(1+s/L).
 \label{eq:rg-load-projection}
\end{equation}
Unused load may be retained to obtain the lower bound.  A record is
\emph{nonwinding} if its load is less than a fixed small constant times
the short period divided by $M_j$; its lift is then injective and its
formula is the bulk formula.  Every other record is called winding,
even when its displayed function is constant or has a short support.
These are the record conventions of
\cite[Section ``Scales, observation, and the class of densities'']{SharpO4},
with the short period used for rectangles.

For a connected graph $X$, put
$1_X=\mathbf1\{r_e\le t_j\text{ for }e\in X\}$.  Local continuations are
taken in fixed sheets on $r_e<4t_j$.  Across a boundary of the chosen
fundamental interval we continue the real $w$ coordinate, rather than
differentiating its discontinuous reduction.  On a small-field stencil,
all local transports agree: a nonflat plaquette has an edge with
$r_e\ge P_j/4$, whereas $P_j/t_j\to\infty$ uniformly along our
trajectories.  Long paths keep their explicitly specified transports.

At anchor $o$, the relative variables are
\[
 y_o(x)=\operatorname{Im}(q_x^{(o)}q_o^{-1}),\qquad
 z_o(x)=\sum_{e\in\text{specified path }o\to x}u_e,
\]
where $q_x^{(o)}$ is transported to the anchor sheet.  Use a fixed basis
of invariant tensors in these variables.  A degree-$n$ polynomial label
has coefficient $a$ and path length
$u=1+\sum_{\nu=1}^n|x_\nu-o|_1$, counting multiplicity, or the comparable
length of its declared paths.  Its coefficient price is
$|a|e^{\sigma u}(1+u)^2$, with a fixed small $\sigma>0$.  The norm is the
supremum over anchors of the sum of these prices.  The ordered slots are
\begin{equation}
 BP_3;\qquad BP_4,I_2;\qquad BP_5,I_3;\qquad
                      BP_6,I_4,B^{-1}J_2.
 \label{eq:R5}
\end{equation}
A subscript denotes homogeneous field degree.  The four groups have
normalized orders one through four: a degree-$n$ term with prefactor
$B^a$ has normalized order $n-2a$, counting each relative field as one
power of $B^{-1/2}$.  Every quadratic orbit
is compensated by directional nearest-edge squares, separately in $q$
and $w$, averaging opposite directions, so that its Hessian vanishes on
every affine field.  Compensators retain their provenance and complete
read sets; their coefficient price uses their own short path.  This is
the directional version of the compensated norm in
\cite[Appendix ``A fixed coefficient norm for the diagonal canonical
map'']{SharpO4}.

A regular error $f_X$ is $C^8$ on the indicated graph domain.  Its local
norm $[f_X]_{8,j}$ is the sum of the supremum norms of derivatives through
order eight, with independent tangent directions satisfying
\[
 |v_x^q|,|v_x^w|\le t_j(1+\operatorname{dist}(x,o))^8.
\]
Quaternion variations use left exponential coordinates; $w$ variations
are ordinary real variations in the fixed sheet.  Its norm and cap are
\[
 \|f\|_{8,j,A}=\sup_o\sum_{X\text{ anchored at }o}
          e^{As_X}[f_X]_{8,j}\le\delta_j:=H_j^{-2.05},\qquad A=64.
\]
Every nonwinding error has zero constant term, zero entire first jet,
and zero Hessian on pairs of affine fields.  Winding errors are exempt
from these normalizations.  In particular it is necessary to remove a
linear $w$ term; quaternion symmetry alone would not remove it.

For the nonsmooth part put $D(q,w,k)=\{e:r_e>t_j\}$.  A covering label
$\ell$ has complete support $P_\ell$, load $s_\ell$, a nonempty inventory
$J_\ell$ of edges, and a bounded measurable weight $k_\ell$.  The weight
is zero unless $J_\ell\subset D$.  Define
\[
 \Xi=\sum_{\substack{\Gamma:\,P_\ell\text{ pairwise disjoint}\\
              \bigsqcup_{\ell\in\Gamma}J_\ell=D}}
                    \prod_{\ell\in\Gamma}k_\ell,
 \qquad
 \sup_x\sum_{\ell:x\in P_\ell}e^{As_\ell}\|k_\ell\|_\infty
                  \le w_j^{\rm cap}:=e^{-p_j^{1/4}}.
\]
Thus every flagged edge is covered exactly once, and $\Xi=1$ when
there are no flags.  The weights may be complex; this is not an optional
polymer gas and need not be positive term by term.

\begin{definition}[Retained density]\label{def:rg-retained-class}
A retained density is an expression
\begin{equation}
 \exp\!\left\{i\theta\sum_{e_t}u_e\right\}
 \exp\{V|\Lambda|-B\mathcal E-X_{\rm can}-X_{\rm err}\}\,\Xi,
 \label{eq:R2}
\end{equation}
with the four kinetic coefficients in the bands above and the numerical
energy $\mathcal E$ formed from the direct terms \eqref{eq:R4} and
the enabled second-stack squares.  If $\mathcal P_a$ denotes a canonical
label, including its displayed prefactor in \eqref{eq:R5}, write
\[
 X_{\rm can}=\sum_a1_{X_a}\mathcal P_a,
 \qquad X_{\rm err}=\sum_X1_Xf_X.
\]
Here $X_a$ is the declared graph of the complete polynomial label,
including its compensators and paths.  Each masked term is defined to
be zero outside its graph domain, so no local continuation is evaluated
there.  The errors and covering sum have the stated norms.  Each polynomial
slot has a fixed coefficient cap, chosen successively in its displayed
order.  The real number $\theta$ is unrestricted.  Bulk scalars and
coefficient extractions are independent of the periods and seam.

All terms obey the listed symmetries, with orbit assignments made on
complete labels.  Reference covering weights vanish whenever a
plaquette containing an edge of their inventory is nonflat.  Brownian
weights restricted to the condition that some such plaquette is
nonflat obey their covering cap with the additional factor $e^{-F_h}$,
where $F_0=0$.  These plaquette tests belong to the complete support.
\end{definition}

For the reference, the covering convention imposes flatness exactly.
Indeed every nonflat plaquette has an edge with $r_e\ge P_j/4>t_j$,
whose mandatory label vanishes.  Elsewhere it makes no change.  The
construction commutes with lattice translation before a seam is fixed.
On a seam rectangle its short formulas use the same seam transports;
moving the seam is a change of coordinates.  Gaussian chart integration
is equivariant under quaternion actions, and the orbit assignments
preserve both the density and its norms.

\subsection{Observation and the phase identity}

When passing from layer $j$ to $j-1$, write
\[
              T=R_j/R_{j-1},\qquad g=B^{-1/2},\qquad
              t=t_j,\qquad t'=t_{j-1}.
\]
This $T$ is a scale conversion, distinct from $T_j$ in
\eqref{eq:rg-scales}.  The propagated output lift in input units is
$\widetilde w'=Tw'$.  Output masks and the threshold in
\eqref{eq:R4} still use $r'_e,t_{j-1}$.  Thus the propagated $w$ square
has coefficient $\chi T^2$, even though its denominator and cutoff are
expressed in output units.

On a flat block transport the inputs into a common lift.  Observe the
$w$ coordinate by a Gaussian of precision $B\chi$ centered at its
$Q$-mean.  For $q$, put $m_Y=Qq(Y)$ and use
\[
 \frac{e^{-B|q'_Y-\widehat m_Y|^2/2}}{z(B)}\,d\omega_3(q'_Y),\qquad
 z(B)=\int_{S^3}e^{-B|q-1|^2/2}d\omega_3(q).
\]
If $|m_Y|\ge1/2$, take $\widehat m_Y=m_Y/|m_Y|$; otherwise sample a
constituent with the $Q$ weights as the target.  The denominator is
independent of the target.  On a block with curvature, choose instead
a constituent minimizing its $Q$-weighted sum of distances to all
constituents in the quotient of $S^3\times\mathbb R$ by
$(q,w)\sim(-q,w+P_j)$.  Use the quotient product metric
\[
 d_{P_j}((q,w),(q^{\prime},w^{\prime}))
 =\min_{n\in\mathbb Z}\left(
 |q-(-1)^nq^{\prime}|^2+|w-w^{\prime}+P_jn|^2\right)^{1/2},
\]
with symmetric randomized ties.  Observe
both its coordinates with the preceding noises.  Such a minimizing
constituent is called a medoid.

Reduce the observed coordinates to the output representative convention,
retaining the shifts.  If two adjacent blocks have flat union, their
common lift induces the coarse-edge integer.  Otherwise choose a coarse
integer giving least quotient distance between the output sites, with
symmetric ties.  The seam sign is included in both rules.  The
construction is a probability kernel.  On a flat input it decouples
exactly into the quaternion observation and a Gaussian observation on
the cover, and preserves both winding integers.

\begin{lemma}[Observation displacement and phase]
\label{lem:rg-observation}
For neighboring blocks $Y,Y'$ and their coarse edge $E$,
\begin{equation}
 r'_E\le\frac C L\sum_{e\subset Y\cup Y'}r_e
       +C\bigl(|\eta_Y|+|\eta_{Y'}|\bigr),
 \label{eq:R8}
\end{equation}
where $\eta$ denotes the observation displacement, with its real
coordinate measured before reduction.  The constant is independent of
large $L$ and of $R_j$.  In conditional integration at fixed outputs
and output integers, the old phase propagates to
$\theta'=\theta LT$.  Its remaining ratio is a product of local
modulus-one factors depending only on sheet integers.  Every nonidentity
factor requires curvature in a block or an adjacent block pair.
\end{lemma}

\begin{proof}
For the quotient metric $d$, put
\[
 A_Y=\sum_{x,z\in Y}Q_Y(x)Q_Y(z)d(s_x,s_z),
 \qquad s_x=(q_x,w_x).
\]
The medoid $m_Y$ satisfies
$\sum_xQ_Y(x)d(m_Y,s_x)\le A_Y$: average the same expression over
candidate constituents.  The triangle inequality therefore gives
\[
 d(m_Y,m_{Y'})\le A_Y+A_{Y'}+
       \sum_{x\in Y,z\in Y'}Q_Y(x)Q_{Y'}(z)d(s_x,s_z).
\]
Bound each distance by a coordinate-path sum of $r_e$.  Smooth bump
weights give total path congestion at most $C/L$ on an edge, proving
the first term in \eqref{eq:R8}.  The argument works before quotienting
on a flat patch.  Normalizing a quaternion mean costs only a fixed
factor when $|m_Y|\ge1/2$; on its fallback branch the weighted path
bound is already bounded below by a positive constant.  Add the two
noise distances to finish the estimate.

We describe the conditional change of variables, since it explains both
the phase and the absence of an extra period factor.  Fix a tree in
each fine block.  Site shifts $n_x$ set all its tree-edge integers to
zero; the remaining integers lie off the tree.  Align the observation
in this lift, using the lifted medoid on a curved block.  Tree equations
fix the differences of the $n_x$, leaving one root shift.  Summation over
observation sheets supplies that shift.  Equivalently, choose the root
representative so that the observed real coordinate lies in the basic
interval; conditioning this coordinate to be $Tw'_Y$ unfolds the root,
and then all tree coordinates, over $\mathbb R$.  Reducing the fine
coordinates reconstructs the original integers and observation sheet,
so this is a bijection.  Its only output scaling Jacobian is $T$ per
block.

On a flat adjacent pair the residual integers inside both blocks vanish,
and all fitted seam integers equal the prescribed coarse integer.
On a curved pair the least-distance condition reads only the quotient
output and its specified integer.  The tree used inside a curved block
cannot enter a flat-pair rule at that block.  Thus tree choices do not
alter the identity; they may be symmetrized identically in both inputs.

Summing temporal real differences in these tree coordinates telescopes.
Since each coarse temporal edge has $L$ crossing fine edges and
$P_j=TP_{j-1}$,
\[
 \sum_{e_t}u_e-LT\sum_{E_t}u'_E
   =P_j\left\{\sum_{\text{internal }e_t}k_e+
       \sum_{E_t}\sum_{e\text{ crossing }E}(k_e-k'_E)\right\}.
\]
Internal residual integers vanish on a flat block, and every crossing
difference vanishes on a flat pair.  Multiplication by $i\theta$ gives
the claimed local phase factors.  They are independent of continuous
field variations in fixed sheets.
\end{proof}

\subsection{The one-step estimates}
\label{subsec:rg-step-estimates}

For two inputs at the same layer and history length, let $\lambda$ be
the vector of differences of their four displayed kinetic coefficients,
Brownian minus reference.  On common label sets, let $U$ be a fixed
triangularly weighted sum of their canonical coefficient distances and
of
\[
 H_j^{-1/2}\delta_j^{-1}\|f_B-f_F\|_{7,j,A},\qquad
 H_j^{-1/2}(w_j^{\rm cap})^{-1}\|k_B-k_F\|_{j,A}.
\]
The covering norm is the support-hit norm just defined.  Each density is
controlled through eight derivatives; comparisons use seven.

\begin{theorem}[Exact step and comparison]\label{thm:rg-step}
Choose the constants in the order specified at the end of this section.
For every sufficiently large $H$, every admitted layer $j\ge1$, every
admitted rectangle or bulk density of Definition~\ref{def:rg-retained-class},
and either seam, the normalized block observation of
Lemma~\ref{lem:rg-observation} has an exact output in the same class,
provided the output kinetic coefficients remain in the next admitted
bands.  In input normalization
for the $w$ coefficients, each of the four kinetic coefficients changes by
\begin{equation}
              \alpha+\kappa/B+O_L(H_j^{-1.05}),
 \label{eq:R6}
\end{equation}
where the direction and field indices are understood.  The bounded
coefficients $\alpha,\kappa$ and the next polynomial coefficients are
given by the full-kernel Taylor map constructed below.

For a Brownian input and a flat reference input on common lists, their
outputs satisfy, for fixed $q_*<1$ and $c>0$,
\begin{equation}
 \begin{split}
 U'&\le q_*U+C_LL^{-h}+C_LH_j^{-c}|\lambda|,\\
 |\lambda'-\lambda|&\le C_LU+C_LL^{-h}
                                      +C_LH_j^{-c}|\lambda|.
 \end{split}
 \label{eq:R7}
\end{equation}
The coefficients in $\lambda'$ use output units, so the second estimate
includes the real-coordinate conversion by $T$.  No discrepancy of the
old phase parameter is required.  The restricted curvature cap renews
with
\begin{equation}
                 F_{h+1}=F_h+c_L\sqrt\beta\,p_j.
 \label{eq:R10}
\end{equation}
Per output site, after the incoming scalar has been propagated, the
scalar increment is
\begin{equation}
 4(L^2-1)\log g+V_L(l_q,l_w,\chi,T,h)+O_L(H_j^{-1}),
 \label{eq:R11}
\end{equation}
where $V_L$ is real, uniformly bounded and Lipschitz on the admitted
bands.  Its arguments may also specify the type of free history.
There is no additional power of $P_j$.
\end{theorem}

The following subsections prove the theorem.  The observation has already
isolated the old phase in curvature regions.  We first obtain a positive
energy reserve for these regions and for all large increments.  This
reserve controls the unbounded real coordinates and integer sums in a
localized Gaussian expansion.  We state the joint product estimate for
that expansion and then verify its hypotheses for the retained class.
Connected summation gives the exact output density; its finite Taylor
part gives the canonical map, while a separate argument contracts the
remaining regular errors.  Finally we restore the low-jet normalizations
and account for winding records and the scalar.

\subsection{The energy reserve near large increments}

We use a local decomposition of the integration domain to keep the large
increments away from the Gaussian Taylor calculation.  Insert a smooth
edge profile equal to one for $r_e\le t/4$ and zero for $r_e\ge t/2$,
and an analogous noise profile at a sufficiently large fixed multiple
of $t$.  A selected complement is an input \emph{seed}; every true
output flag is a mandatory seed.  The list of selected complements is
a fixed primary pattern during each integral.  A configuration produced
by subsequently opening other profiles does not change that pattern.
Join seeds whose $50M_j$ neighborhoods meet.  For each connected
component freeze all fine blocks meeting its $3M_j$ neighborhood.  The
remaining coordinates are the exterior Gaussian coordinates.  A
\emph{core} is the factor assigned to one such component, together with
its frozen integrations and retained tests.

Here is the precise geometric prescription we import from
\cite[Appendix ``Local sector factors and their complete supports'',
Subsections ``One metric and one assignment of each factor'' and
``The read sets of the prepared factors'']{SharpO4}.  Assign numerical
energy rows within $5M_j$ of the seeds to the core, stationary squares
within $8M_j$, and retained default profiles within $14M_j$.  The
stationary squares are the residual squares from the Gaussian completion
below.  Select Gaussian affine rows farther than $2.5M_j$ from the seeds.  Include every
positive and negative test in the factor that uses it.  Covariance-window
reads include the entire window and the tests deciding which sites in
it are exterior; determinant compensations include the corresponding
hole neighborhoods.  A noncore factor retains a $20M_j$ enlargement of
its arguments and joining record.  Long old paths and determinant chains
retain all their endpoints and windows.  Local composition radii are
chosen so their sum is less than $M_j/100$; the fixed outer margins then
contain every read.  Freezing whole blocks adds only $O(L)$ to these
radii.  We enlarge both hole tests and determinant neighborhoods by that
same margin, and take $H$ large after $L$.  This is a finite change of
the recorded support dilation $D_*$.

The energy estimate behind this prescription is algebraic.  Let $X$ be
the input stack applied to $(\nabla q,u)$ before masking, $X'$ the output
stack applied to $(\nabla q',Tu')$, and $Y$ the observation displacement
with its observation-square weights.  All $w$ entries carry
$\sqrt\chi$.  Let $\zeta$ have direct output rows clipped by
$\min(1,t'/r'_E)$ and the enabled second-stack output rows.  For the
numerically truncated operators put $(\nu,\omega)=\mathcal M_s\zeta$.
Then
\[
                |\zeta|+|\omega|\le Ct,
                      \qquad |\nu|\le Ct/L.
\]
These bounds apply to each row of the stacks, uniformly in the volume.
The subscript $s$ denotes numerical stencils.  With a row's direct and
second entries understood, the input energy plus observation energy
minus output energy is the sum of
\begin{align}
 \mathcal F_e&=S_t(e)+\tfrac12m_e|X_{2,e}|^2
                         -\nu_e\cdot X_e+\tfrac12|\nu_e|^2,
 \nonumber\\
 \mathcal B_Y&=\tfrac12|Y_Y-\omega_Y|^2,
 \nonumber\\
 \mathcal C_E&=\zeta_E\cdot X'_E-\tfrac12|\zeta_E|^2
                         -S_{t'}(E)-\tfrac12m'_E|X'_{2,E}|^2,
 \nonumber\\
 \mathcal U_E&=\zeta_E\cdot[\mathcal M_s^*(X,Y)-X']_E,
 \nonumber\\
 \mathcal N_i^{\rm row}&=\tfrac12| (\mathcal N_s\zeta)_i|^2,
 \qquad
 \mathcal Z_E=\tfrac12\zeta_E\cdot
       [(I-\mathcal M_s^*\mathcal M_s-\mathcal N_s^*\mathcal N_s)
                                                    \zeta]_E.
 \label{eq:rg-ledger}
\end{align}
Here $m_e,m'_E$ are the second-stack masks.  The cancellation of cross
terms separates the direct, observation, and output reserves in
$\mathcal F,\mathcal B,\mathcal C$ from the defects in
$\mathcal U,\mathcal Z$; the $\mathcal N^{\rm row}$ term is nonnegative.
The identity uses actual transposes and products of the chosen matrices, so
\eqref{eq:rg-ledger} holds even when the transport has holonomy.
Transport observation rows through their block-tree lift and transport
other rows through their declared paths.  Their absolute exponential
envelopes are unchanged by signs.  Free differential identities are
used only when the entire stencil, including intervening blocks and
pair constraints, has a common flat lift.

\begin{lemma}[Reserve for flagged and curved regions]
\label{lem:rg-reserve}
After the affine Gaussian squares have been selected and subtracted,
each core with $s_c$ seeds retains energy at least $c_Lt^2s_c$,
as in \cite[Lemma ``Core reserve'']{SharpO4}, and, in addition, an unspent
fixed fraction of
\begin{equation}
 ct\sum_{\text{input flags}}r_e
       +ct'\sum_{\text{output flags}}r'_E
       +c\sum_{\text{noise seeds}}|\eta_Y|^2.
 \label{eq:R9}
\end{equation}
The energy is multiplied by $B$ in the density.  The estimate holds on
the retained tests and on the closed supports of the profile derivatives
used below.  If the core contains curvature it retains at least
$c tP_j$ before multiplication by $B$.
\end{lemma}

\begin{proof}
The direct $\mathcal F$ row pays $ct^2$ at an edge seed and $ctr_e$ at a
true input flag.  Its cross term has coefficient $Ct/L$.  A failed
second-stack mask can be charged to one witness $r_e>t$.  A witness
hits at most $C(1+\log(r_e/t))^2\le C'r_e/t$ rows; the extra $D_0$
exclusion only changes this constant.  Exponential column sums therefore
bound the total loss by
\[
                         \frac{Ct}{L}\sum_{r_e>t}r_e.
\]
Large observation displacements have their square reserve, after the
fixed noise threshold is increased.  The direct $\mathcal C$ row pays
$ct'r'_E$ on an output flag and is zero otherwise; its second-stack
part cancels.

On a flat stencil, the free identities give the quaternion defect bound
\[
 C_Lt e^{-c_LM_j}
       +Ct\sum_Y e^{-c\operatorname{dist}(E,Y)}|m_Y-\widehat m_Y|.
\]
The real-coordinate identity is linear and exact before truncation.
After truncation, subtract a common real constant on the enlarged
stencil.  Differences of coordinate values cost at most its length
times the largest increment or noise there.  The error is therefore
\[
 C_LM_j^Dt e^{-c_LM_j}
       \bigl(1+\max_{\text{enlarged stencil}}(r_e,|\eta_Y|)\bigr).
\]
This estimates coefficients of finite stencils, not remote unbounded
values.  A maximum above the default threshold lies in the same core
and is paid by its direct or noise reserve.  For the quaternion mean,
the block path estimate gives $(C/L)\sum r_e$; on blocks with chords
below $T_j$ its error improves to $C_LT_j^2$.  Every other block has a
witness above $T_j$, to which the $Ct/L$ charge is assigned.

It remains to estimate a stencil that sees curvature, where a
free differential identity is unavailable.  Its $\mathcal Z$ row costs
$C_Lt^2$.  For $\mathcal U$, use the absolute coefficient envelopes in
\eqref{eq:rg-ledger}.  The fine part of $\mathcal M_s$ has row sum
$C/L$, and the input stack has bounded column sums.  Summing curved
rows consequently gives a fine-input charge
\[
                  \frac{Ct}{L}\sum_e a_er_e,
\]
where the nonnegative incidence weights have bounded column
multiplicity.  The observation and coarse-stack column sums similarly
give $Ct\sum_E a'_Er'_E+Ct\sum_Ya_Y|\eta_Y|$.  Substituting
\eqref{eq:R8}, each neighboring block pair uses an input edge only a
bounded number of times, so the coarse charge again gives
$(Ct/L)\sum r_e$ plus noise terms.  The constants multiplying $L^{-1}$
are independent of $L$.  Large increments are paid by the direct rows;
large noises are paid by their squares.  All remaining increments,
noises, and $\mathcal Z$ rows cost at most $C_LM_j^Dt^2$ per curvature
witness.  Such a witness has $r_e\ge P_j/4$.  The ratio of this last
cost to $tP_j$ is $C_LM_j^Dt/P_j$, which tends to zero.  Output curvature
also forces fine curvature in the enlarged stencil: the pair-lift rules
carry a flat neighborhood to a flat neighborhood, including across a
four-block corner.

Choose $L$ to spend less than a fixed fraction of the direct reserve
on all $L^{-1}$ charges.  Only flat small-angle stencils are used for
selected affine squares; their subtraction costs
$C_LM_j^Dt^3$ per row.  Frozen angular extensions are cut off at
$C_Lt$.  Exterior principal logarithms retain the chart and alignment
protectors, with alignment imposed also on the real coordinate.
The remaining losses are smaller than the reserve once
\[
 tM_j^D\ll1,
 \qquad M_j^D(T_j^2/t+e^{-c_LM_j}/t)\ll1.
\]
These choices leave \eqref{eq:R9} and a fixed $ctP_j$ whenever there is
curvature.  On a no-output-flag extension through $r'_E<4t'$, the
least-distance output integer, if queried, is unique; a flat-pair
constraint is a fixed-sheet constraint.  The same estimates thus hold
on the required derivative supports.
\end{proof}

The reserve is the point at which the extra coordinate and its integer
curvature enter the nonlinear argument.  We next use it to construct
absolutely convergent local expansions.  This step concerns complete
products under one Gaussian law; estimates on individual expectations
would not suffice.

\subsection{Localized Gaussian integration and a joint product estimate}

On exterior sites write $q_x=e^{\xi_x}q'_{[x]}$ in the small principal
logarithm and use ordinary relative coordinates for $w_x$.  Retained
defaults give alignment within $C_Lt$.  They also fix all residual
integers on any exterior block or pair remote from seeds.  Perform this
elimination on pairs adjoining the frozen boundary as well, retaining
the whole pair in the margin.  The remaining integers and hard
observation predicates are local reads.  The core keeps their complete
halo, including contrary output-status answers.  Off that halo every
formula uses fitted sheet coordinates, rather than untested integer
sums.

For the globally extended selected-row formula, smoothly clip frozen
real relative coordinates and output graph differences beyond a larger
$C_Lt$ plateau, as well as clipping frozen quaternion angles.  These
cutoffs equal the original relative arguments on the retained selected-row
supports.  They are part of the background $F$ in the affine rows
$D\xi+F$ below and make the extended mean uniformly bounded in all
frozen arguments.
Combine the three quaternion tangent coordinates and the real
relative coordinate into one four-component variable $\xi$.
The selected affine rows have the form $D\xi+F$.  For quaternion
coordinates rotate a row into the nearby output frame; transport signs
multiply its background too, so the tangent matrix $D$ is the ordinary
free matrix.  For the real coordinate it is already linear.  Thus
$A_{\mathrm{fl}}=D^*D$ is a principal restriction of the numerical precision
$d^*\mathsf A_s^*\mathsf A_sd+Q^*Q$, direct-summed with the real
coordinate and its $\chi$ weights.  Every row incident on an exterior
column is selected.  Block Poincar\'e, including $Q$, gives
\[
                 c_LI\le A_{\mathrm{fl}}\le C_LI
\]
uniformly in the seed pattern and periods.  This matrix depends on the
geometry, not the frozen fields or output values.

For an exterior site $x$, let $\Omega_x$ be its exterior sites inside
the numerical inverse window, and define
\[
 (\Pi_u)_{yx}=\mathbf1_{y\in\Omega_x}
               (A_{{\mathrm{fl}},\Omega_x}+u)^{-1}_{yx},\qquad
 E_u=(A_{\mathrm{fl}}+u)\Pi_u-I.
\]
The window argument of
\cite[Section ``The energy identity and the massive Gaussian'',
window bounds and ratio-series identities]{SharpO4} gives exponentially
weighted row and column estimates
\[
 \|(A_{\mathrm{fl}}+u)^{-1}\|+\|\Pi_u\|\le C_L(1+u)^{-1},\qquad
 \|E_u\|\le C_Le^{-c_LM_j}(1+u)^{-1}.
\]
Set
\[
 C_s=(\Pi_0+\Pi_0^*)/2,\qquad
 \mu=-\Pi_0^*D^*F,\qquad r_s=D\mu+F.
\]
For large $H$, $C_s$ and all its principal marginals have upper and
positive lower bounds at fixed $L$.  Substitute $\xi=\mu+gZ$ under the
single Gaussian probability law $Z\sim N(0,C_s)$.  The exact Gaussian
ratio left in the integrand is
\[
 \exp\!\left\{\tfrac12Z^*(C_s^{-1}-A_{\mathrm{fl}})Z
       -g^{-1}(A_{\mathrm{fl}}\mu+D^*F)^*Z\right\}.
\]
The scalar determinant is expanded by the resolvent identity and the
series for $\log(I+E)$; every nonlocal correction retains its literal
window path and has at least one residual $E$ hop.  With any fixed
support exponent, the sum of envelopes of paths with $r$ hops is
\[
 C_L\operatorname{poly}(M_j,r)e^{CAr}
                          (C_Le^{-c_LM_j})^r.
\]
The resolvent integral uses the integrable bound $(1+u)^{-2}$.  The
empty-pattern leading determinant is extracted as a bulk scalar on
all sites; its missing frozen-site values and local hole corrections
are assigned to the core.

The prediction and stationary estimates needed in the factor bounds are
\begin{equation}
 \begin{split}
 |\mu|&\le C_LM_j^Dt,\\
 |d q^{\rm pred}|+|d w^{\rm pred}|&\le Ct/L+C_LM_j^Dt^2,\\
 |r_s|&\le C_LM_j^D(t^2+te^{-c_LM_j})
 \end{split}
 \label{eq:rg-prediction-local}
\end{equation}
pointwise on remote flat stencils; the first bound holds for the global
extended formula.  To see the remote estimates, replace the window and numerical
kernels by the full empty-pattern kernels.  The first jet of the
prediction is exactly $\mathsf P_h$ by
\eqref{eq:rg-positive-completion}; all remaining quaternion terms are
quadratic.  Equation~\eqref{eq:rg-free-prediction} supplies the first
and second fine-difference gains.  For $w$, subtract its common constant
before estimating.  Differences of all free data and these formulas
cost
\[
                        C_L(L^{-h}+|\lambda|/H_j).
\]

We record precisely which part of the prepared-factor machinery is
imported.  This also identifies the estimates required for the new real
coordinate.

\begin{lemma}[Joint integration contract]\label{lem:rg-joint}
Fix a primary pattern and the complete-support convention above.  Form
local factors from: nonlinear rows minus their affine squares;
stationary squares; old polynomial and regular-error terms; exceptional
numerical defects; determinant and Gaussian-ratio chains; Haar Jacobian
ratios; opened default-profile failures; and old covering weights.
Include the compulsory cores and their frozen integrations.

Suppose their ordinary smooth factors have envelopes
$C_Lg^r\operatorname{poly}(M_j,p_j,s)$ of positive order $r$; ratio
exponents have bounds
\[
 |Q_\alpha(Z)|\le e_\alpha
        \left(1+\sum_{x\in I_\alpha}|Z_x|^2\right);
\]
their eighth-root support-hit sums are sufficiently small; and every selection of $n$ distinct
failure tests has simultaneous probability at most
$\exp(-c_Lp_j^2n/M_j^D)$.  Suppose further that the frozen integral of each core satisfies
\[
 \int|B_c|\le\exp\{-c_Lp_j^2s_c+
 C_LM_j^D(1+|\log g|)s_c\}
\]
uniformly in the Gaussian arguments,
while retaining the output-flag and curvature reserves.  Here $B_c$ is
the core integrand; its integral includes the frozen quaternion and real
variables and the residual integer and tag sums, with Gaussian and output
arguments held fixed.  Suppose also that differentiation through
hard dependencies can be represented by one simultaneous transport
whose fixed-order $L^4$ scores cost a polynomial in total load and a
fixed power of $g^{-1}$.  Whenever such a negative power occurs, retain
a core, failure, or exponentially accurate factor to pay it.
Assume explicitly that every allowed differentiated product is a sum of
polynomially many terms retaining the ordinary and ratio envelopes and
the closed failure supports, apart from these score factors; a marked
discrepancy uses its replacement envelope.

Then every finite selection has a joint integral bound equal to the
product of its stated majorants times a fixed polynomial in total load.
This holds through eight normalized output derivatives for empty output
inventory, and through seven derivatives for one marked input or
parameter discrepancy.  With nonempty output inventory it gives the
sup norm and the corresponding discrepancy bound.  A core with $s_c$ seeds retains
$e^{-c_Lp_j^2s_c}$ after fixed losses, and each failure test retains
$e^{-c_Lp_j^2/M_j^D}$; ratio chains may use their eighth-root envelopes.  Disjoint complete supports factor
exactly; overlapping factors require no independence.
\end{lemma}

\begin{proof}
The abstract Gaussian estimate is
\cite[Appendix ``Joint majorants for Gaussian prepared factors'',
Lemma ``A joint Gaussian product estimate'' and Corollary ``Fixed-order
derivatives and an exceptional reserve'']{SharpO4}.  We recall its short
argument.  If $|Q_\alpha|\le e_\alpha U_\alpha$ and $U_\alpha\ge1$, then
\[
 |e^{Q_\alpha}-1|\le\sqrt{e_\alpha}
                  e^{2\sqrt{e_\alpha}U_\alpha}.
\]
When $0\le C_s\le c_*I$ and
$\sup_x\sum_{\alpha:x\in I_\alpha}\sqrt{e_\alpha}
<(64c_*)^{-1}$, the Gaussian determinant identity controls the fourth
moment of the product of these envelopes.  The individual smallness
condition
\[
 (2+4c_*|I_\alpha|)\sqrt{e_\alpha}
                  \le\tfrac14\log(e_\alpha^{-1})
\]
gives $\|\prod_\alpha(e^{Q_\alpha}-1)\|_4
\le\prod_\alpha e_\alpha^{1/4}$.  H\"older with exponents $4,4,2$
combines the score, ratio product, and simultaneous-failure indicator.
A fixed number of derivatives has polynomially many placements in
total load.  Pay its negative power of $g$ from one retained reserve;
weakening the other reserves uniformly gives the asserted bounds.
For a discrepancy, telescope into one marked factor and use its
replacement envelope after factoring out the discrepancy size.  No
bound relative to the old factor is used; that factor may vanish.

The eighth-root sums follow with their support prices, not merely from
an unweighted chain sum.  The number of literal paths with $r$ numerical hops is at most
\[
 \operatorname{poly}(M_j,r)(C_LM_j^D)^{r+O(1)}.
\]
H\"older on this finite path set and the chain bound with eight times
the requested exponent give
\[
 \sum_{\alpha:\,r\text{ hops}}e^{As_\alpha}e_\alpha^{1/8}
 \le\operatorname{poly}(M_j,r)
       e^{-c_LM_jr/8+O_L(r\log M_j)+CAr},
\]
which is summable and small for large $H$.

Exact factorization for disjoint complete supports is
\cite[Appendix ``Local sector factors and their complete supports'',
Lemma ``Local marginals and factorization'']{SharpO4}.  Its hypotheses
are precisely the inclusion of all window and hole queries and all
positive and contrary status reads.  They make the frozen sets disjoint,
the retained marginal covariance block diagonal, and the window data
unchanged when the other factor is omitted.  This is an assertion about
these marginal laws, not conditional independence under an unspecified
Gaussian law.
\end{proof}

\subsection{Verification of the joint product bounds}

We now apply Lemma~\ref{lem:rg-joint} to the factors produced by our block
integration.  Its hypotheses require four checks: an integrable core
bound, positive-order smooth and ratio envelopes, simultaneous failure
tails, and derivative bounds under one common transport.  The first check
is essential for the real lift: a small pointwise core weight alone
cannot be integrated over its unbounded coordinates.
Each old inventory edge is marked in its selected input core and must
be supplied exactly once.  At every other selected input edge retain
the test $r_e\le t$; nonprimary edges already have $r_e\le t/2$.
These rules reproduce the mandatory input cover.  The core contains its
assigned rows with affine squares subtracted, nearby stationary
squares, all primary and nearby default tests, frozen normalization
reciprocals, and leading-window determinant compensations.  It retains
every output-status answer needed by those formulas.  Attached smooth
protectors keep every exterior quaternion argument inside its injective
chart and are never opened separately.

Every frozen block retains its observation-noise anchor, and every
remaining unfitted off-tree integer lies in owned rows; the $3M_j$
freeze and $5M_j$ ownership margins ensure both assertions.  A fixed
fraction of \eqref{eq:R9}, after multiplication by $B$, leaves
\[
 \exp\!\left\{-a_0\left(
    \sum_{\text{tree edges}}|\Delta w|+
    \sum_{\text{residual integer reads}}|\Delta w+P_jk|
       +\sum_{\text{frozen blocks}}|\eta_Y^w|\right)\right\}
\]
for fixed $a_0>0$.  Inserting these factors at default-size values costs
at most $e^{C_LM_j^Ds_c}$ for a core with $s_c$ seeds.  Since
$P_j\ge P_{\min}>0$,
\[
 \sup_{a\in\mathbb R}\sum_{k\in\mathbb Z}
             e^{-a_0|a+P_jk|}
 \le 1+\frac{2}{1-e^{-a_0P_{\min}}}.
\]
Sum off-tree integers with this bound.  Tree differences and one common
translate per block then give real coordinates with unit Jacobian.
The observation displacement anchors that translate: a $Q$-mean changes
by the translate because its weights sum to one, and a fixed medoid
changes by the same amount.  In the medoid case first sum the bounds
over its at most $L^2$ possible sites, without imposing the medoid
test.  The frozen integral is consequently bounded, uniformly in all
remaining Gaussian and output arguments, by
\begin{equation}
       \exp\{-c_Lp_j^2s_c+C_LM_j^D(1+|\log g|)s_c\}.
 \label{eq:rg-frozen-entropy}
\end{equation}
It retains the output-flag and curvature reserves of
\eqref{eq:R9}, with a smaller fixed constant.  There is no power of
$R_j$ in this estimate.

For the remaining factors use the following explicit preparations.
A remote term means that its complete projected graph is farther than
$14M_j$ from seeds.
\begin{enumerate}[label=\textup{(\roman*)}]
\item On a smooth unassigned row use the actual row minus its selected
      affine square, with a smooth cutoff equal to one on the retained
      sector and supported at angles $C_Lt$.  Multiplied by $B$, its
      cubic remainder has envelope
      $C_Lg\operatorname{poly}(M_j,p_j)$.  Mean vertices have the same
      order.  Stationary squares and quaternion Haar Jacobians start at
      order $g^2$.  The real coordinate has Lebesgue Jacobian one.
\item A polynomial of degree $n$ and prefactor $B^a$ on a path of
      length $u$ has envelope
      $C_L|a_X|g^{n-2a}\operatorname{poly}(M_j,p_j,u)$, because its
      relative variables are at most $Ctu$.  In particular $BP_3$ has
      order one.  Old errors use their stated derivative norms with
      fixed logarithmic costs.  Smooth remote masks are supported inside
      their old graph extensions through $2t$; a nonremote hard mask
      remains attached to the compulsory or exceptional factor that
      required it.
\item Numerical defects in selected rows, full-kernel coefficient
      replacements, and the $\mathcal N^{\rm row},\mathcal Z$
      corrections have envelopes $C_Lg^{-D}e^{-c_LM_j}$, with their
      complete paths and support prices.  There are no omitted tails
      of the actual energy: it was defined by numerical rows.  Full
      kernels occur only in Taylor coefficients or stencils with
      bounded normalized increments.
\item Determinant and Gaussian-ratio chains retain their literal paths
      and the ratio envelopes in Lemma~\ref{lem:rg-joint}.
\item Open the remaining edge and noise defaults, and exterior chart
      profiles, as one plus a signed failure factor.  The remote mean
      uses its smooth normalized branch with a larger plateau.  By
      \eqref{eq:rg-prediction-local}, failure and its derivative supports
      force $|Z|\ge c_Lp_j$ on the associated stencil.  Selecting a separated subcollection of the queried
      stencils and using the bounded covariance gives
      $\exp(-c_Lp_j^2n_d/M_j^D)$ for $n_d$ failures.  The real-coordinate
      defaults satisfy the identical Gaussian implication.
\item Old covering weights keep their sup envelopes, complete
      inventories and protectors; their measurable functions are not
      directly differentiated.
\end{enumerate}
Each preparation agrees with the original product on the full sector.
The empty-pattern determinant, Gaussian powers, Haar constant, and
observation denominators have already been extracted, so no ordinary
factor of order zero remains.

For hard dependencies we use the protected-coordinate transport of
\cite[Section ``Prepared factors and a joint integral estimate'',
transport identity, and Appendix ``Joint majorants for Gaussian
prepared factors'']{SharpO4}.  Its content is a change of variables in
one Gaussian marginal on the union of all read coordinates, chosen to
hold all measurable physical-field arguments fixed.  The inverse
quaternion chart is smooth on the attached protector.  For $w$, add the
linear transport keeping the fine real coordinate fixed; its inverse
differential costs at most $C_L/g$.  Tree integers, medoid tags, and
flat-pair sheet constraints are held fixed.  The least-integer output
tests are not differentiated when the inventory is nonempty, and on a
no-flag extension they are locally constant through $r'<4t'$.
Consequently a fixed number of derivatives produces Gaussian scores
whose $L^4$ bounds are a polynomial in total load and a fixed power of
$g^{-1}$.  The window difference identity gives the same statement for
parameter transport.  Ellipticity alone is not being used as a score
estimate; the bounded chart inverses, extended means, and differentiated
window kernels provide its hypotheses.

The protected transport leaves frozen variables fixed.  Its coefficients
are uniformly bounded by the relative-coordinate cutoffs just specified;
parameter derivatives introduce polynomials in physical increments and
noise displacements, never bare root coordinates or bare integers.
For the frozen integrations, take absolute values after this transport,
use the smooth, ratio, and failure envelopes, and integrate the frozen
variables using \eqref{eq:rg-frozen-entropy} before applying the Gaussian
product estimate.  Distinct compulsory footprints have disjoint frozen
variables, and their bounds are pointwise in $Z$.  Powers of increments
or noise introduced by a fixed number of derivatives are paid by the
remaining exponential reserve.  One may first restrict all frozen
integrals and integer sums and then pass to the unrestricted integral
by this domination.  In particular, no product of overlapping optional
expectations has been replaced by the product of expectations.

We next verify comparisons, including the old phase.  Interpolate core
exponents at the \emph{same physical arguments}, keeping frozen real
variables on all of $\mathbb R$ and quaternion arguments on their
protectors.  Both endpoint core estimates hold there and on closed
profile derivative supports.  The differentiated exponent has factor
$C_L(L^{-h}+|\lambda|/H_j)$ times fixed polynomial costs in increments
and noises, which the reserve pays before frozen integration.  For
changing Gaussian maps use a common marginal transport; interpolate
numerical means linearly and ratio exponents with their exceptional
envelopes.  An interpolated stack identity is unnecessary.  The leading
predicted increments interpolate linearly, with quaternion composition
error $C_LM_j^Dt^2$, so remote failure implications persist.  Smooth
vertex formulas have common cutoff geometry and can be interpolated
directly.  These verifications give the single-discrepancy version of
Lemma~\ref{lem:rg-joint} with fixed polynomial costs, not a product of
costs growing with the number of factors.

For the old phase, Lemma~\ref{lem:rg-observation} puts every nontrivial
ratio in a curved input patch.  The retained halo contains an input
flag of that plaquette and the old covering factor supplying it.  The
corresponding reference factor is zero.  We may therefore replace the
old reference phase in that compulsory product by the Brownian phase
without changing the reference integrand.  The difference is still
charged to the marked old covering discrepancy.  This remains true on
empty-output derivative supports: the zero is a retained sheet and
inventory condition.  In the remote exterior, where some defaults have
been opened, fitted sheets give phase ratio one.  Thus no old $\theta$
or its discrepancy enters the norm.

The same argument propagates the curvature restriction.  If an output
inventory edge meets a nonflat plaquette, its compulsory component and
full halo force fine curvature, since pair lifts preserve a flat
neighborhood through four-block corners.  That component retains an
old covering factor restricted to its curvature condition and the
additional reserve $Bc_LtP_j$.  An output-empty background factor or a
forbidden background subtraction is never substituted for this
compulsory component.  We have now verified every hypothesis of
Lemma~\ref{lem:rg-joint} for the lifted, complex-valued integration.

\subsection{Connected sums and the finite Taylor map}

For an optional hard exponential with attached protector and alignment
profiles $\chi_{\rm att}$, use exactly
$1+\chi_{\rm att}(e^V-1)$.  The attached profiles belong to the
selected letter; they are not first multiplied into $e^V$ and then
opened, which would create an unpriced order-zero failure.
Open the other prepared exponentials as $1+(e^V-1)$ and join factors
whose complete read supports intersect.  The preceding joint bound
majorizes these connected products.  We use the rooted-tree summation
of \cite[Section ``Connected sums and exact reassembly'', tree
condition and exact background-exclusion identity]{SharpO4} in the
following form.  If support-hit majorants have a sufficiently small
sum after their exponential load prices, the sum over connected
selections anchored at a fixed site converges absolutely; intersecting
supports cost $C_LM_j^2$ times load.  Projection of old records obeys
\eqref{eq:rg-load-projection}.  The unused input exponential therefore
pays the change of support exponent, every fixed load power, and the
finite enlargement of read sets.  Core positioning, halo and status
choices cost $\exp(C_Ls_c\log M_j)$, paid by
\eqref{eq:rg-frozen-entropy}.  Marking an old factor, including a
curvature-restricted one, costs only a load polynomial.

For clarity, the exact reassembly distinguishes components with and
without output inventory.  Take the logarithm of the latter to form the
regular output exponent.  In the presence of a compulsory output
component, a disallowed background component contributes its negative
logarithmic term; open that term and attach it to the compulsory
component.  The remaining components form a mandatory cover whose
inventories partition exactly the output flagged set.  This is an
identity of absolutely convergent series, so complex coefficients cause
no change.  It neither divides by a restricted covering sum nor assumes
its positivity.

Inflate an ordinary factor of normalized order $r$ by $g^{-.96r}$
before applying the tree condition.  Its remaining
$g^{.04r}\operatorname{poly}(M_j,p_j)$ still tends to zero faster than
the fixed overlap price.  Terms of total order at least five therefore
retain $g^{4.8}$ before the final fixed losses.  Exceptional accuracies,
core reserves and shared tails pay their derivative, entropy and tree
prices.  The exact expansion is thereby controlled at all numbers of
factors, while only a finite Taylor polynomial is required for the
coefficient map.

\begin{lemma}[Canonical map and contraction]\label{lem:rg-canonical}
The principal Taylor map through normalized order four produces exactly
the slots \eqref{eq:R5}, scalar terms, an imaginary temporal linear
$w$ term, and the four affine kinetic corrections in \eqref{eq:R6}.
At their own order and displayed prefactor, the diagonal response of a
degree-$n\ge3$ slot has norm at most $CL^{2-n}$; a compensated quadratic
has norm at most $C/L$.  The constants are independent of $L$ before
it is chosen.  The map is triangular in the order \eqref{eq:R5}, and
its paired coefficient differences obey a strict contraction plus
$C_L(L^{-h}+|\lambda|/H_j)$.
\end{lemma}

\begin{proof}
The principal terms involve at most four origins.  Strengthen their
masks to one common ball and use an even fluctuation guard.  The shared
Gaussian tail pays the complement.  Count each relative output field
and each $gZ$ as one power of $g$, and each kinetic prefactor $B$ as
$g^{-2}$.  The vertices are the nonlinear mean, actual row minus affine
square, stationary square, old polynomial, and Haar Jacobian vertices
of \cite[Section ``The canonical coefficients'']{SharpO4}, with the
real coordinate in direct sum and the additional input cubic vertex.
Taylor's theorem through normalized order four gives remainder
$C_Lg^5\operatorname{poly}(M_j,p_j)$, including normalized derivatives
and paired estimates.  Restoring full Gaussian polynomial moments and
full kernels costs exponential accuracy in $M_j$.  Connected Wick
coefficients converge in the same exponential path norm: each covariance
connection and prediction carries the exponential moments of
Lemma~\ref{lem:rg-free}, and all paths are retained.

A connected term with $v$ prefactored vertices has at least $v-1$
intervertex contractions.  Each contraction costs $g^2$; hence at most
one $g^{-2}$ prefactor remains.  At normalized order $r\le4$, the
maximal field degree is $r+2$.  The possibilities are therefore
$BP_3$ at order one, $BP_4,I_2$ at order two, $BP_5,I_3$ at order
three, and $BP_6,I_4,B^{-1}J_2$ at order four, besides scalar and linear
terms.  Linear $w$ terms occur at orders one and three.  Spatial
reflection leaves only the temporal direction, and conjugation makes
their coefficients imaginary.  The affine quadratic symmetry gives
exactly the four Hessians already displayed.

Remove those Hessians at orders two and four.  Their coefficients are
$\alpha$ and $\kappa/B$ in \eqref{eq:R6}.  A removal changes the shape
at which the next free kernel is evaluated.  Accordingly subtract the
Taylor expansion of that changed kinetic term through order four: the
quartic chart-degree response times $\alpha$, and, at quadratic degree
and order $B^{-1}$, the second parameter derivative applied to the
square of the kick, are both included.  The aggregate affine Hessian
is linear in the four coefficients, so these corrections remove exactly
the intended Hessian.  Rewrite slot prefactors at the new $B$, including
the lower-order prefactor change of $BP_3$, and convert real-coordinate
powers by $T$.  This prescription defines the finite Taylor map without
an implicit choice of coupling convention.  Finally compensate each
remaining quadratic orbit by its directional nearest-edge tags.

For the diagonal response within the same slot and at its displayed
prefactor, an old coefficient transfers only by linear prediction.
A relative coordinate over displacement $r$ becomes
\[
 \sum_z[\mathsf P_h(o+r,z)-\mathsf P_h(o,z)]\,y_O(z),
                  \qquad O=[o],
\]
and the same formula holds for the real variable.  The first difference
bound gives $C|r|/L$ with the spare exponential path weight.  A degree-$n$
label therefore gains $L^{-n}$; summing the $L^2$ fine anchors gives
$CL^{2-n}$.  For a quadratic packet, expand each difference into its
linear affine part and a second-difference remainder.  The product of
the two affine parts cancels before absolute values are taken, by the
packet's affine normalization.  The remainder costs $L^{-3}$, hence
$C/L$ after anchor summation.  The cost of output compensation is bounded
by the path's second moment.  This is the proof of
\cite[Appendix ``A fixed coefficient norm for the diagonal canonical
map'', Lemma ``One norm and constants chosen before the block
size'']{SharpO4}, with directional tags in place of a single isotropic
tag.  Its hypotheses are precisely
\eqref{eq:rg-free-prediction} and affine cancellation, which hold here.

The factors $T^n$ are near one.  Contractions lower degree and can feed
only later slots at a given order; nonlinear products have larger
normalized order.  Shape and prefactor subtractions respect this
triangular order, because prediction preserves affine fields and hence
the diagonal quadratic response has zero bulk affine Hessian.  All
nondiagonal bounds are finite for fixed $L$.  Choose caps successively,
including the finite initial caps from Section~\ref{sec:initialization},
and then triangular discrepancy weights.  Taking $L$ large gives the
strict diagonal contraction; free-history and parameter differences
have the stated factor by Lemma~\ref{lem:rg-free}.
\end{proof}

\subsection{Contraction of the regular errors}

It remains to control the functional error, rather than only its Taylor
coefficients.  Remove first the isolated linear transfer of each old
regular error.  The hypotheses of
\cite[Lemma ``Contraction of old regular errors'', in Section ``An exact
renormalization map'']{SharpO4} are zero constant and entire first jet,
zero affine Hessian, and the prediction estimates.  They hold in our
fixed local lifts.  We give the relevant derivative argument, since
this is where an unremoved linear real-coordinate term would fail.

For a label of load $s\le L^{1/4}$ the output graph contains its entire
local ball.  Subtract an anchor quaternion frame and a common real
lift.  Let $A$ be the predicted relative displacement and
$\mathfrak b$ its centered affine approximation.  In old normalized
direction units,
\[
                    \|A\|\le C/L+o(1),\qquad
          \|A-\mathfrak b\|\le C/L^2+o(1).
\]
For distance $r\le L$, divide the first and second difference estimates
$Ct'r/L,Ct'r^2/L^2$ by $t(1+r)^8$.  For $r>L$, the bounded row of
$\mathsf P_h$ gives the stronger weight ratio
$(1+r/L)^8/(1+r)^8\le CL^{-8}$.  The affine term has the same bound.
Nonlinear chart terms have an extra $t$ with a fixed power of $M_j$;
there is no such nonlinearity in $w$.

On an even Gaussian guard, omitting fluctuations through order two
costs
\[
 C_LM_j^D(1+s)^Dp_j^{-2}[f_X]_{8,j}.
\]
The odd term vanishes, and each fluctuation insertion supplies a factor
$1/p_j$ in the old direction norm.  Let $\mathcal B_X$ be the Hessian of the relative
error at zero.  The relative error itself has zero value and first
derivative there.  Write $A_i$ and $A_{12}$ for the first and mixed
second derivatives of the prediction $A$ in normalized output directions.
The value and first two output derivatives of the error have leading terms
\[
 \tfrac12\mathcal B_X(A,A),\qquad
 \mathcal B_X(A,A_i),\qquad
 \mathcal B_X(A_1,A_2)+\mathcal B_X(A,A_{12}).
\]
Replacing every argument by the corresponding derivative of
$\mathfrak b$ gives zero, since those fields remain affine.  Each
remaining bilinear has one $C/L^2+o(1)$ factor and one $C/L+o(1)$
factor; the cubic Taylor remainder has the same gain.  For derivatives
of orders three through eight, retain the fluctuation and apply the
chain rule directly.  There are either at least three first map
variations or a higher map variation, giving $C/L^3+o(1)$ without using
more old derivatives than the requested order.  The $L^2$ anchor sum
gives $q_1=C/L+o(1)$.  For longer or winding labels use direct composition
and the unused old support exponential; its $e^{-cL^{1/4}}$ pays the
anchor count and every load polynomial.  Local real-coordinate
variations have the bounded first row independent of the chosen sheet.

Let $R$ be the norm of the raw output regular error through eight
derivatives, before its low jets are removed, and let $R^\Delta$ be
the corresponding paired difference norm through seven derivatives.
Together with the high-order connected sums, the preceding contraction
argument gives
\begin{equation}
 \begin{split}
 R&\le q_1\delta_j+C_Lg^{4.6},\\
 R^\Delta&\le q_1\delta_j^\Delta+
       C_L\operatorname{poly}(M_j,p_j)\delta_j
                     (L^{-h}+|\lambda|/H_j)+C_Lg^{4.5}\tau.
 \end{split}
 \label{eq:rg-raw-remainders}
\end{equation}
Here $\delta_j^\Delta$ is the undivided error distance, and $\tau$ is
the canonical distance plus the relative error and covering distances
\emph{without} the factor $H_j^{-1/2}$, together with
$L^{-h}+|\lambda|/H_j$.  A nonisolated old error has its stronger actual
cap and at least one additional positive-order factor, so lies in the
fifth-order remainder.  Shape changes in the kinetic term are expanded
analytically on its fixed stencil masks; the chart-jet truncation error
is exponential in $M_j$.  None of these operations differentiates a
moving flag predicate.

\subsection{Normalization, winding records, and closure}

The raw error need not yet have zero low jets.  We restore that
normalization without changing the density, while keeping the
unbounded real coordinate out of bounded covering corrections.
Strengthening a regular mask has the exact correction supported on its
complement; every removed term then meets an output flag and is
reassembled into its mandatory covering label by the connected-sum
identity above.

For a nonwinding raw error at anchor $o$, first subtract its value and
its entire linear real-coordinate polynomial.  Write the latter using
real increments on declared paths from $o$ to each argument.  The paths
may be chosen with length at most $C(1+|x-o|)$ after a fixed load
enlargement.  To construct them, follow the old connected cube record
and add a shortest shortcut from the anchor whenever the current graph
distance exceeds three times the ambient distance.  Each added length
is at most a fixed multiple of the travel since the preceding shortcut,
by the triangle inequality.  Add connections inside cubes and a full
enlarged anchor cube for nearby arguments.  The total record length
remains $O(M_js)$; orbit symmetrization uses its common graph.  If the
enlargement is no longer injective on the torus, retain it as a winding
record, as described below.

If $a_x$ are the linear coefficients, the directional norm gives
\begin{equation}
              \sum_x t'(1+|x-o|)^8|a_x|
                            \le C[f_X]_{8,j-1}.
 \label{eq:rg-linear-removal}
\end{equation}
The norm of a path sum depends only on its two endpoints.  Hence this
removal costs a fixed factor in the raw norm.  For mask-complement
corrections use odd clipped increments, equal to the true increment
through $4t'$ and bounded by $Ct'$ outside a larger interval.  Use their
clipped squares for affine tags.  Quaternion nearest-edge tags may be
clipped with the same plateau.  Common clipped tags preserve the exact
cancellation of any orbit whose total tagged Hessian is zero.  These
choices make mask-complement corrections small even when the tag
coefficient is of order zero.

Apply the same clipping to canonical linear terms being extracted as
phase.  By translation and reflection, their bulk sum is a constant
times the sum of the clipped temporal real increments.  Replace it by
the corresponding true-increment phase, assigning the difference to the
label that owns each flagged edge.  The difference exponent is
imaginary, so this multiplier has modulus one.  Only the newly extracted
linear coefficient is compared or differentiated; the old propagated
phase is unchanged.

Next subtract the four affine Hessians.  The norm of a unit affine
direction is at most $C(t')^{-1}$ in the derivative convention, so the
coupling correction is bounded by $C(t')^{-2}$ times the raw norm.
Insert the change of the numerical kinetic energy at the shifted
coefficients.  Allocate its nearest-edge compensators to each path in
proportion to that path's directional affine Hessians.  Their sums are
exactly the four required coefficients.  On the common path mask the
remaining expression has zero affine jet and costs at most a fixed
factor of the raw norm, independent of $L$: the free row moments and
their shape derivatives are uniform.  This is the kinetic reset of
\cite[Section ``Normalization, periods, and closure'', kinetic-reset
identity]{SharpO4}, with four directional tags.

There is one additional tail term in this reset.  Changing the direct
energy \eqref{eq:R4} produces an unbounded exponent outside its mask.
Attach this exponent only to the label owning that flag.  A coefficient
change of size $C_L$ costs at most $C_Lt'r'_E$, and
\eqref{eq:R9} has the factor $B$ needed to absorb it.  Active
second-stack rows and their parameter derivatives are bounded by $Ct'$,
so their mask complements have the usual small tree bound.  Comparing
these multipliers introduces only fixed powers of flagged increments,
again paid by the reserve.  For phase multipliers the same argument
uses the newly extracted coefficient discrepancy; it never uses an
old phase coefficient.

All scalar, phase and Hessian extractions so far use bulk lists.  A
connected calculation whose complete record is injective on the torus
agrees with the corresponding bulk calculation, including either seam
transport.  A disagreement requires output load at least
\[
                c_*\min_i(m_iL^{j-1})/M_{j-1}.
\]
A reserved part of the support exponential pays that record.  Retain
the disagreement as an actual winding regular error or covering term,
including any period-dependent constant, linear term, or Hessian.
Use clipped tags and include their edges in the winding error's mask;
on its extension through $4t'$ the unclipped chart formula is then
valid.  There is no period-dependent phase extraction.  Already winding
inputs retain enough winding load under
\eqref{eq:rg-load-projection}.  Thus the bulk scalar and coefficient
maps are common to every admitted period and both seams.

We can now read off the bounds in Theorem~\ref{thm:rg-step}.  Output
covering components have norm at most $\exp(-p_j^{3/10})$ before division
by their cap.  In the curvature restriction they also contain
\[
                       e^{-F_h-c_LBt_jP_j}.
\]
Since $B\asymp H_j$ and $P_j=\pi\sqrt{\beta/H_j}$,
$Bt_jP_j\asymp\sqrt\beta\,p_j$, proving \eqref{eq:R10}.  The estimate
applies to the sum of restricted labels, with one retained old
curvature factor and one fresh reserve.  It asserts no independent
factor for every output bad edge or every plaquette.  The reference-zero
rule follows from the same compulsory-product argument.

Divide \eqref{eq:rg-raw-remainders} by the output error cap.  Up to
comparable $B,H_j$ constants,
\[
           g^{4.6}/\delta_j=O(g^{.5}),\qquad
           g^{4.5}/\delta_j=O(g^{.4}).
\]
The additional $H_j^{-1/2}$ weight in $U$ absorbs the fixed logarithmic
loss on a free-history discrepancy.  In the $\tau$ term its conversion
back to the unweighted error distance cancels the output weight and
leaves the positive $g^{.4}$ gain.  Affine extraction acts on the
undivided error and costs $(t')^{-2}$, so its individual correction is
$O_L(H_j^{-1.05})$ and its paired correction has the form allowed in
\eqref{eq:R7}.  Covering differences have stronger suppression than
their cap.  The canonical contraction is
Lemma~\ref{lem:rg-canonical}.  Since $|T-1|=O_L(H_j^{-1})$, the change
from input to output normalization also has the permitted
$H_j^{-c}|\lambda|$ size.  These facts prove \eqref{eq:R6} and
\eqref{eq:R7}.

For completeness, here is the order of constants closing all estimates.
First fix the direct coefficients and the deletion-radius constants
needed in the microscopic construction, including $D_0$.  Reserve
finitely many larger support exponents for intermediate tree, mask,
orbit and winding operations, and fix their geometric dilation factors.
Then choose $L$ large for the free operator bounds, the $L^{-1}$ reserve
charges, the diagonal contractions, and the support-exponent projection.
Choose the slot caps, including the finite initialization caps, and the
triangular discrepancy weights.  Let $D$ exceed all fixed powers of
$M_j$, load, score variables and normalization entropy used above.
Choose $P_0>20$ sufficiently large in terms of $D$, and finally take $H$
large.  In particular, uniformly for $H_j\ge H$,
\begin{align}
 &C_Le^{-c_LM_j}\operatorname{poly}(M_j)\ll1,
 &&C_Lt_jM_j^D\ll1,\nonumber\\
 &C_LM_j^D/p_j^2\ll1,
 &&C_LM_j^D(T_j^2/t_j+e^{-c_LM_j}/t_j+t_j)\ll1,
 \nonumber\\
 &C_LM_j^D(1+\log g^{-1})<c_Lp_j^2/4,
 &&c_Lp_j^2/M_j^D\gg p_j^{3/10}+D\log g^{-1}+D\log M_j,
 \nonumber\\
 &C_Lg^{.04}\operatorname{poly}(M_j,p_j)\ll M_j^{-2},
 &&C_Lg^{4.6}/\delta_j\ll1.
 \label{eq:rg-finite-choices}
\end{align}
They are the finite-choice inequalities of
\cite[Section ``Normalization, periods, and closure'']{SharpO4}.
Each follows from a positive power of $H_j$ dominating fixed logarithms,
or from taking $P_0$ above the fixed powers of $M_j$.  The extra
curvature inequality $C_LM_j^Dt_j/P_j\ll1$ follows uniformly from
$H_j\le C\beta$ and the same final increase of $H$.  All $L$-independent
constants used to choose $L$ were identified before this step.  Enlarging
$P_0$ does not enlarge the fixed exponent $D$.

Finally consider the scalar.  Gaussian integration over the aligned
exterior contributes four tangent coordinates per fine site, and the
normalized observation removes four coordinates per output site.  This
gives $4(L^2-1)\log g$.  The real output conversion contributes only its
bounded factor $T$.  The empty-pattern determinant has the full-kernel
value up to exponential accuracy; the window determinant estimates and
uniform shape analyticity make its remainder $V_L$ bounded and
Lipschitz.  Scalar Taylor corrections begin at normalized order two,
and the remaining scalar errors are controlled by
\eqref{eq:rg-raw-remainders}.  Conjugation symmetry makes the bulk scalar
real.  The observation-sheet bijection preserves the single global
circle mode; it introduces no factor $P_j$ per block.  This proves
\eqref{eq:R11} and completes the proof of Theorem~\ref{thm:rg-step}.
